% Source: https://artofproblemsolving.com/community/c3h1927802_2012_aime_i_problem_6
% Stable identity: aops-topic-1927802-q01
\textbf{Why $z^{142}=1$.} Substituting, $z=w^{11}=\left(z^{13}\right)^{11}=z^{143}$. If $z=0$, its imaginary part would be $0$; but $0<\dfrac{m\pi}{n}<\pi$, so $\sin\dfrac{m\pi}{n}>0$. Hence $z\ne0$, and dividing by $z$ gives $z^{142}=1$.

Then $|z|^{142}=1$, so $|z|=1$ and $z=\cos\theta+i\sin\theta$. By De Moivre's theorem $z^{142}=\cos142\theta+i\sin142\theta=1$, so $142\theta$ is a multiple of $2\pi$:
\[
\theta=\dfrac{2k\pi}{142}=\dfrac{k\pi}{71}\quad(k=0,1,\dots,141),\qquad \operatorname{Im}z=\sin\dfrac{k\pi}{71}.
\]

\textbf{Finding $n$.} The imaginary part is positive, so $0<k<71$. Two angles in $(0,\pi)$ with equal sines are equal or supplementary, so
\[
\dfrac {m}{n}=\dfrac{k}{71}\quad\text{or}\quad\dfrac {m}{n}=\dfrac{71-k}{71}.
\]
Both numerators lie between $1$ and $70$ and $71$ is prime, so both fractions are already in lowest terms. Hence $n=\boxed{71}$.

Such $z,w$ exist: $z=\cos\dfrac{\pi}{71}+i\sin\dfrac{\pi}{71}$ and $w=z^{13}$ give $w^{11}=z^{143}=z^{142}z=z$.
